\IfFileExists{stacks-project.cls}{%
\documentclass{stacks-project}
}{%
\documentclass{amsart}
}

% For dealing with references we use the comment environment
\usepackage{verbatim}
\newenvironment{reference}{\comment}{\endcomment}
%\newenvironment{reference}{}{}
\newenvironment{slogan}{\comment}{\endcomment}
\newenvironment{history}{\comment}{\endcomment}

% For commutative diagrams we use Xy-pic
\usepackage[all]{xy}

% We use 2cell for 2-commutative diagrams.
\xyoption{2cell}
\UseAllTwocells

% We use multicol for the list of chapters between chapters
\usepackage{multicol}

% This is generally recommended for better output
\usepackage{lmodern}
\usepackage[T1]{fontenc}

% For cross-file-references
\usepackage{xr-hyper}

% Package for hypertext links:
\usepackage{hyperref}

% For any local file, say "hello.tex" you want to link to please
% use \externaldocument[hello-]{hello}
\externaldocument[introduction-]{introduction}
\externaldocument[conventions-]{conventions}
\externaldocument[sets-]{sets}
\externaldocument[categories-]{categories}
\externaldocument[topology-]{topology}
\externaldocument[sheaves-]{sheaves}
\externaldocument[sites-]{sites}
\externaldocument[stacks-]{stacks}
\externaldocument[fields-]{fields}
\externaldocument[algebra-]{algebra}
\externaldocument[brauer-]{brauer}
\externaldocument[homology-]{homology}
\externaldocument[derived-]{derived}
\externaldocument[simplicial-]{simplicial}
\externaldocument[more-algebra-]{more-algebra}
\externaldocument[smoothing-]{smoothing}
\externaldocument[modules-]{modules}
\externaldocument[sites-modules-]{sites-modules}
\externaldocument[injectives-]{injectives}
\externaldocument[cohomology-]{cohomology}
\externaldocument[sites-cohomology-]{sites-cohomology}
\externaldocument[dga-]{dga}
\externaldocument[dpa-]{dpa}
\externaldocument[sdga-]{sdga}
\externaldocument[hypercovering-]{hypercovering}
\externaldocument[schemes-]{schemes}
\externaldocument[constructions-]{constructions}
\externaldocument[properties-]{properties}
\externaldocument[morphisms-]{morphisms}
\externaldocument[coherent-]{coherent}
\externaldocument[divisors-]{divisors}
\externaldocument[limits-]{limits}
\externaldocument[varieties-]{varieties}
\externaldocument[topologies-]{topologies}
\externaldocument[descent-]{descent}
\externaldocument[perfect-]{perfect}
\externaldocument[more-morphisms-]{more-morphisms}
\externaldocument[flat-]{flat}
\externaldocument[groupoids-]{groupoids}
\externaldocument[more-groupoids-]{more-groupoids}
\externaldocument[etale-]{etale}
\externaldocument[chow-]{chow}
\externaldocument[intersection-]{intersection}
\externaldocument[pic-]{pic}
\externaldocument[weil-]{weil}
\externaldocument[adequate-]{adequate}
\externaldocument[dualizing-]{dualizing}
\externaldocument[duality-]{duality}
\externaldocument[discriminant-]{discriminant}
\externaldocument[derham-]{derham}
\externaldocument[local-cohomology-]{local-cohomology}
\externaldocument[algebraization-]{algebraization}
\externaldocument[curves-]{curves}
\externaldocument[resolve-]{resolve}
\externaldocument[models-]{models}
\externaldocument[functors-]{functors}
\externaldocument[equiv-]{equiv}
\externaldocument[pione-]{pione}
\externaldocument[etale-cohomology-]{etale-cohomology}
\externaldocument[proetale-]{proetale}
\externaldocument[relative-cycles-]{relative-cycles}
\externaldocument[more-etale-]{more-etale}
\externaldocument[trace-]{trace}
\externaldocument[crystalline-]{crystalline}
\externaldocument[spaces-]{spaces}
\externaldocument[spaces-properties-]{spaces-properties}
\externaldocument[spaces-morphisms-]{spaces-morphisms}
\externaldocument[decent-spaces-]{decent-spaces}
\externaldocument[spaces-cohomology-]{spaces-cohomology}
\externaldocument[spaces-limits-]{spaces-limits}
\externaldocument[spaces-divisors-]{spaces-divisors}
\externaldocument[spaces-over-fields-]{spaces-over-fields}
\externaldocument[spaces-topologies-]{spaces-topologies}
\externaldocument[spaces-descent-]{spaces-descent}
\externaldocument[spaces-perfect-]{spaces-perfect}
\externaldocument[spaces-more-morphisms-]{spaces-more-morphisms}
\externaldocument[spaces-flat-]{spaces-flat}
\externaldocument[spaces-groupoids-]{spaces-groupoids}
\externaldocument[spaces-more-groupoids-]{spaces-more-groupoids}
\externaldocument[bootstrap-]{bootstrap}
\externaldocument[spaces-pushouts-]{spaces-pushouts}
\externaldocument[spaces-chow-]{spaces-chow}
\externaldocument[groupoids-quotients-]{groupoids-quotients}
\externaldocument[spaces-more-cohomology-]{spaces-more-cohomology}
\externaldocument[spaces-simplicial-]{spaces-simplicial}
\externaldocument[spaces-duality-]{spaces-duality}
\externaldocument[formal-spaces-]{formal-spaces}
\externaldocument[restricted-]{restricted}
\externaldocument[spaces-resolve-]{spaces-resolve}
\externaldocument[formal-defos-]{formal-defos}
\externaldocument[defos-]{defos}
\externaldocument[cotangent-]{cotangent}
\externaldocument[examples-defos-]{examples-defos}
\externaldocument[algebraic-]{algebraic}
\externaldocument[examples-stacks-]{examples-stacks}
\externaldocument[stacks-sheaves-]{stacks-sheaves}
\externaldocument[criteria-]{criteria}
\externaldocument[artin-]{artin}
\externaldocument[quot-]{quot}
\externaldocument[stacks-properties-]{stacks-properties}
\externaldocument[stacks-morphisms-]{stacks-morphisms}
\externaldocument[stacks-limits-]{stacks-limits}
\externaldocument[stacks-cohomology-]{stacks-cohomology}
\externaldocument[stacks-perfect-]{stacks-perfect}
\externaldocument[stacks-introduction-]{stacks-introduction}
\externaldocument[stacks-more-morphisms-]{stacks-more-morphisms}
\externaldocument[stacks-geometry-]{stacks-geometry}
\externaldocument[moduli-]{moduli}
\externaldocument[moduli-curves-]{moduli-curves}
\externaldocument[examples-]{examples}
\externaldocument[exercises-]{exercises}
\externaldocument[guide-]{guide}
\externaldocument[desirables-]{desirables}
\externaldocument[coding-]{coding}
\externaldocument[obsolete-]{obsolete}
\externaldocument[fdl-]{fdl}
\externaldocument[index-]{index}

% Theorem environments.
%
\theoremstyle{plain}
\newtheorem{theorem}[subsection]{Theorem}
\newtheorem{proposition}[subsection]{Proposition}
\newtheorem{lemma}[subsection]{Lemma}

\theoremstyle{definition}
\newtheorem{definition}[subsection]{Definition}
\newtheorem{example}[subsection]{Example}
\newtheorem{exercise}[subsection]{Exercise}
\newtheorem{situation}[subsection]{Situation}

\theoremstyle{remark}
\newtheorem{remark}[subsection]{Remark}
\newtheorem{remarks}[subsection]{Remarks}

\numberwithin{equation}{subsection}

% Macros
%
\def\lim{\mathop{\mathrm{lim}}\nolimits}
\def\colim{\mathop{\mathrm{colim}}\nolimits}
\def\Spec{\mathop{\mathrm{Spec}}}
\def\Hom{\mathop{\mathrm{Hom}}\nolimits}
\def\Ext{\mathop{\mathrm{Ext}}\nolimits}
\def\SheafHom{\mathop{\mathcal{H}\!\mathit{om}}\nolimits}
\def\SheafExt{\mathop{\mathcal{E}\!\mathit{xt}}\nolimits}
\def\Sch{\mathit{Sch}}
\def\Mor{\mathop{\mathrm{Mor}}\nolimits}
\def\Ob{\mathop{\mathrm{Ob}}\nolimits}
\def\Sh{\mathop{\mathit{Sh}}\nolimits}
\def\NL{\mathop{N\!L}\nolimits}
\def\CH{\mathop{\mathrm{CH}}\nolimits}
\def\proetale{{pro\text{-}\acute{e}tale}}
\def\etale{{\acute{e}tale}}
\def\QCoh{\mathit{QCoh}}
\def\Ker{\mathop{\mathrm{Ker}}}
\def\Im{\mathop{\mathrm{Im}}}
\def\Coker{\mathop{\mathrm{Coker}}}
\def\Coim{\mathop{\mathrm{Coim}}}

% Boxtimes
%
\DeclareMathSymbol{\boxtimes}{\mathbin}{AMSa}{"02}

%
% Macros for moduli stacks/spaces
%
\def\QCohstack{\mathcal{QC}\!\mathit{oh}}
\def\Cohstack{\mathcal{C}\!\mathit{oh}}
\def\Spacesstack{\mathcal{S}\!\mathit{paces}}
\def\Quotfunctor{\mathrm{Quot}}
\def\Hilbfunctor{\mathrm{Hilb}}
\def\Curvesstack{\mathcal{C}\!\mathit{urves}}
\def\Polarizedstack{\mathcal{P}\!\mathit{olarized}}
\def\Complexesstack{\mathcal{C}\!\mathit{omplexes}}
% \Pic is the operator that assigns to X its picard group, usage \Pic(X)
% \Picardstack_{X/B} denotes the Picard stack of X over B
% \Picardfunctor_{X/B} denotes the Picard functor of X over B
\def\Pic{\mathop{\mathrm{Pic}}\nolimits}
\def\Picardstack{\mathcal{P}\!\mathit{ic}}
\def\Picardfunctor{\mathrm{Pic}}
\def\Deformationcategory{\mathcal{D}\!\mathit{ef}}
\usepackage{cleveref}
\crefname{originalreference}{original reference}{original references}
\Crefname{originalreference}{Original reference}{Original references}
\makeatletter
\newlabel{algebraic-definition-relative-representable-property}{{03YK}{source}{}{sourcepointer-1}{}}
\newlabel{algebraic-definition-relative-representable-property@cref}{{[originalreference][0][]03YK}{[1][0][]source}{}{}{}}
\makeatother
\crefname{originalreference}{original reference}{original references}
\Crefname{originalreference}{Original reference}{Original references}
\makeatletter
\newlabel{spaces-limits-lemma-surjection-is-enough}{{0CM6}{source}{}{sourcepointer-2}{}}
\newlabel{spaces-limits-lemma-surjection-is-enough@cref}{{[originalreference][0][]0CM6}{[1][0][]source}{}{}{}}
\makeatother
\crefname{originalreference}{original reference}{original references}
\Crefname{originalreference}{Original reference}{Original references}
\makeatletter
\newlabel{spaces-limits-proposition-characterize-locally-finite-presentation}{{04AK}{source}{}{sourcepointer-3}{}}
\newlabel{spaces-limits-proposition-characterize-locally-finite-presentation@cref}{{[originalreference][0][]04AK}{[1][0][]source}{}{}{}}
\makeatother
\crefname{originalreference}{original reference}{original references}
\Crefname{originalreference}{Original reference}{Original references}
\makeatletter
\newlabel{spaces-limits-remark-limit-preserving}{{05N0}{source}{}{sourcepointer-4}{}}
\newlabel{spaces-limits-remark-limit-preserving@cref}{{[originalreference][0][]05N0}{[1][0][]source}{}{}{}}
\makeatother
\title{For representable morphisms, limit preservation on objects equals local finite presentation}
\author{The Stacks Project Authors}
\date{stacks-project:criteria.tex:local-snapshot-20261012; GNU FDL 1.2 or later}
\begin{document}
\maketitle
\noindent\textit{Editorial note.} No omission markers in-span; proof complete as written.\par\medskip
\noindent\textit{Allowed standard prerequisites (editorial).}\begin{itemize}
\item Characterization of locally finite presentation. spaces-limits-proposition-characterize-locally-finite-presentation (cross-chapter, named). Descends morphisms along directed limits of affines in the (2)=>(1) direction.
\item Surjection criterion. spaces-limits-lemma-surjection-is-enough (cross-chapter, named). Yields local finite presentation from surjectivity of the colimit map in the (1)=>(2) direction.
\end{itemize}\medskip
\ifdefined\hypertarget\hypertarget{sourcepointer-1}{}\fi
\noindent\textit{Source reference.} Original reference 03YK: focused-seeds/geometry/S9-crit-representable-by-spaces-limit-preserving/ref-tags.txt, source lines 5--5. Navigation pointer only: the counted proof cites this statement by \textbackslash{}ref as a standing input; the deduction the counted proof has to make is carried out inside the counted span, and no part of that deduction is delegated to this pointer.\par\medskip
\ifdefined\hypertarget\hypertarget{sourcepointer-2}{}\fi
\noindent\textit{Source reference.} Original reference 0CM6: focused-seeds/geometry/S9-crit-representable-by-spaces-limit-preserving/ref-tags.txt, source lines 6--6. Navigation pointer only: the counted proof cites this statement by \textbackslash{}ref as a standing input; the deduction the counted proof has to make is carried out inside the counted span, and no part of that deduction is delegated to this pointer.\par\medskip
\ifdefined\hypertarget\hypertarget{sourcepointer-3}{}\fi
\noindent\textit{Source reference.} Original reference 04AK: focused-seeds/geometry/S9-crit-representable-by-spaces-limit-preserving/ref-tags.txt, source lines 7--7. Navigation pointer only: the counted proof cites this statement by \textbackslash{}ref as a standing input; the deduction the counted proof has to make is carried out inside the counted span, and no part of that deduction is delegated to this pointer.\par\medskip
\ifdefined\hypertarget\hypertarget{sourcepointer-4}{}\fi
\noindent\textit{Source reference.} Original reference 05N0: focused-seeds/geometry/S9-crit-representable-by-spaces-limit-preserving/ref-tags.txt, source lines 8--8. Navigation pointer only: the counted proof cites this statement by \textbackslash{}ref as a standing input; the deduction the counted proof has to make is carried out inside the counted span, and no part of that deduction is delegated to this pointer.\par\medskip
\begin{lemma}
\label{lemma-representable-by-spaces-limit-preserving}
Let $p : \mathcal{X} \to \mathcal{Y}$ be a $1$-morphism of categories
fibred in groupoids over $(\Sch/S)_{fppf}$. If $p$ is
representable by algebraic spaces, then the following are equivalent:
\begin{enumerate}
\item $p$ is limit preserving on objects, and
\item $p$ is locally of finite presentation (see
Algebraic Stacks,
Definition \ref{algebraic-definition-relative-representable-property}).
\end{enumerate}
\end{lemma}

\begin{proof}
Assume (2). Let $U = \lim_{i \in I} U_i$ be the directed limit
of affine schemes $U_i$ over $S$, let $y_i$ be an object of $\mathcal{Y}$
over $U_i$ for some $i$, let $x$ be an object of $\mathcal{X}$ over $U$,
and let $\gamma : p(x) \to y_i|_U$ be an isomorphism. Let
$X_{y_i}$ denote an algebraic space over $U_i$ representing the $2$-fibre
product
$$
(\Sch/U_i)_{fppf} \times_{y_i, \mathcal{Y}, p} \mathcal{X}.
$$
Note that $\xi = (U, U \to U_i, x, \gamma^{-1})$ defines an object of
this $2$-fibre product over $U$. Via the $2$-Yoneda lemma $\xi$ corresponds
to a morphism $f_\xi : U \to X_{y_i}$ over $U_i$. By
Limits of Spaces, Proposition
\ref{spaces-limits-proposition-characterize-locally-finite-presentation}
there exists an $i' \geq i$ and a morphism $f_{i'} : U_{i'} \to X_{y_i}$
such that $f_\xi$ is the composition of $f_{i'}$ and the projection
morphism $U \to U_{i'}$. Also, the $2$-Yoneda lemma tells us that
$f_{i'}$ corresponds to an object
$\xi_{i'} = (U_{i'}, U_{i'} \to U_i, x_{i'}, \alpha)$ of
the displayed $2$-fibre product over $U_{i'}$ whose restriction to
$U$ recovers $\xi$. In particular we obtain an isomorphism
$\gamma : x_{i'}|U \to x$. Note that $\alpha : y_i|_{U_{i'}} \to p(x_{i'})$.
Hence we see that taking $x_{i'}$, the isomorphism
$\gamma : x_{i'}|U \to x$, and the isomorphism
$\beta = \alpha^{-1} : p(x_{i'}) \to y_i|_{U_{i'}}$
is a solution to the problem.

\medskip\noindent
Assume (1). Choose a scheme $T$ and a $1$-morphism
$y : (\Sch/T)_{fppf} \to \mathcal{Y}$. Let
$X_y$ be an algebraic space over $T$ representing the $2$-fibre product
$(\Sch/T)_{fppf} \times_{y, \mathcal{Y}, p} \mathcal{X}$.
We have to show that $X_y \to T$ is locally of finite presentation.
To do this we will use the criterion in
Limits of Spaces, Remark \ref{spaces-limits-remark-limit-preserving}.
Consider an affine scheme $U = \lim_{i \in I} U_i$ written as the
directed limit of affine schemes over $T$.
Pick any $i \in I$ and set $y_i = y|_{U_i}$. Also denote $i'$ an element
of $I$ which is bigger than or equal to $i$. By the $2$-Yoneda lemma
morphisms $U \to X_y$ over $T$ correspond bijectively
to isomorphism classes of pairs $(x, \alpha)$ where $x$ is an object
of $\mathcal{X}$ over $U$ and $\alpha : y|_U \to p(x)$ is an isomorphism.
Of course giving $\alpha$ is, up to an inverse, the same thing as giving
an isomorphism $\gamma : p(x) \to y_i|_U$.
Similarly for morphisms $U_{i'} \to X_y$ over $T$. Hence (1) guarantees
that the canonical map
$$
\colim_{i' \geq i} X_y(U_{i'}) \longrightarrow X_y(U)
$$
is surjective in this situation. It follows from
Limits of Spaces, Lemma \ref{spaces-limits-lemma-surjection-is-enough}
that $X_y \to T$ is locally of finite presentation.
\end{proof}

\end{document}
